\documentclass[11pt]{article}
\usepackage[a4paper,margin=1.8cm]{geometry}
\usepackage[T1]{fontenc}
\usepackage{lmodern}
\usepackage{amsmath,amssymb}
\usepackage{booktabs,tabularx,array}
\usepackage{float,algorithm,algpseudocode}
\usepackage{xcolor}
\usepackage[hidelinks]{hyperref}
\usepackage{fancyhdr}
\definecolor{ink}{HTML}{17354A}
\floatstyle{ruled}
\restylefloat{algorithm}
\pagestyle{fancy}
\fancyhf{}
\fancyhead[L]{\small\color{ink}FlashAttention on the accelerator}
\fancyhead[R]{\small Corrected pseudocode}
\fancyfoot[C]{\small\thepage}
\setlength{\headheight}{14pt}
\setlength{\parindent}{0pt}
\setlength{\parskip}{5pt}
\setlength{\tabcolsep}{5pt}
\renewcommand{\arraystretch}{1.16}
\newcommand{\I}[1]{\operatorname{int#1}}
\newcommand{\F}{\operatorname{fp32}}
\newcommand{\clip}{\operatorname{clip}}
\newcommand{\round}{\operatorname{round}}
\newcommand{\rmax}{\operatorname{rowmax}}
\newcommand{\rsum}{\operatorname{rowsum}}
\newcommand{\VSync}{\textsc{VSync}}
\newcommand{\SafeDot}{\textsc{SafeDot}}
\algrenewcommand\algorithmiccomment[1]{\hfill\textit{\footnotesize // #1}}
\begin{document}

{\Large\bfseries\color{ink}FlashAttention forward pass}\par
{\large Corrected prefill pseudocode for the supplied accelerator}

The kernel keeps K/V in SRAM, visits query tiles in order, and updates a running
softmax result. The control CPU schedules DMA, both systolic arrays, and vector
routines. The tile sizes remain $B_r=128$, $B_c=256$, with $d=128$.

\textbf{Scope and arithmetic contract.}
Q/K/V projections and RoPE finish before this kernel. Inputs occupy 16 bits in
DRAM; their format and decoding scales are supplied by the caller. The vector
CPU decodes them and requantizes to symmetric signed INT8 with zero point zero.
This adds approximation relative to 16-bit attention. Assume the arrays accept
signed INT8 and return exact signed INT16 results when representable, and the
vector CPU supports FP32 and INT32 arithmetic. The bit-width table alone does
not establish these arithmetic semantics.

For decoded values $X$, use a scalar scale and deterministic rounding:
\[
s_X=\begin{cases}\max|X|/127,&\max|X|>0,\\1,&\text{otherwise},\end{cases}
\qquad X_8=\I{8}\!\left(\clip(\round(X/s_X),-127,127)\right).
\]
K and V each use one scale per resident K/V head; Q uses one scale per query
tile. Output scale $s_O>0$ is supplied per query head. All running softmax
statistics and the unnormalized output $U$ use FP32.

\textbf{Scheduling unit.}
Algorithm 1 processes one batch item and one K/V head, reusing that head for its
four query heads. The control CPU invokes it for all 16 batch items and eight
K/V heads, then signals kernel completion. Head-major DRAM storage makes every
query/output tile contiguous. Other layouts require explicit transfer/packing
changes. Attention scores and probabilities are never written to DRAM.

\begin{tabularx}{\linewidth}{@{}p{3.25cm}X@{}}
\toprule
Operation & Meaning on this architecture \\
\midrule
\textsc{DmaRead/Write} & CPU waits for DMA acceptance, configures a contiguous copy,
starts it, and receives a completion event. Configuration costs 200 cycles per copy. \\
\VSync$(f,\ldots)$ & CPU launches vector routine $f$ and waits for completion.
Communication costs 300 cycles per launch. \\
\textsc{SaDot} & CPU waits for array acceptance and launches one physical dot-product
tile. Communication costs 100 cycles per launch. All pointers refer to SRAM. \\
\textsc{Wait/WaitAll} & CPU waits until the named operations and their memory writes
finish. Separate engines may overlap; no hardware command queue is assumed. \\
\bottomrule
\end{tabularx}

\textbf{SRAM allocation for $N=2048$, one resident K/V head.}
\begin{tabularx}{\linewidth}{@{}Xr@{}}
\toprule
Buffers & Reserved size \\
\midrule
Original K and V, 16-bit & 1 MiB \\
Quantized $K_8$ and packed $V_8^\top$ & 0.5 MiB \\
Two signed digit planes for each quantized K/V operand & 1 MiB \\
Query/output ping-pong buffers, score/probability tiles, INT32 products,
FP32 state, digit packing, INT16 array partials, events and alignment & 2 MiB \\
\midrule
Total reservation & 4.5 MiB $<16$ MiB \\
\bottomrule
\end{tabularx}
The workspace includes a 256 KiB raw-partial buffer, enough for either call to
\SafeDot. The remaining SRAM is available for additional prefetching. For other
sequence lengths, pad K/V to a multiple of 256 and recompute the allocation.

\clearpage
\begin{algorithm}[H]
\caption{Control-CPU program for one batch item and one shared K/V head}
\label{alg:main}
\footnotesize
\begin{algorithmic}[1]
\Require Four query heads $Q_h\in\mathbb R^{N\times d}$, shared $K,V\in\mathbb R^{N\times d}$,
output addresses $O_h$, scales $s_{O,h}$; $d=128$, $B_r=128$, $B_c=256$.
\State $T_r\gets\lceil N/B_r\rceil$; allocate the SRAM buffers from page 1.
\State $e_K\gets\Call{DmaRead}{K_{\rm DRAM},K_{16},2Nd}$.
\State $e_V\gets\Call{DmaRead}{V_{\rm DRAM},V_{16},2Nd}$; $\Call{WaitAll}{e_K,e_V}$.
\State $(K^{\rm dig},V^{\rm dig},s_K,s_V)\gets\VSync(\textsc{PrepareKV},K_{16},V_{16})$.
\Statex \hspace{\algorithmicindent}\textit{PrepareKV decodes, quantizes, zero-pads, and packs reusable K/V digit operands.}
\For{$h=0,1,2,3$}
    \State $w[0],w[1]\gets\textsc{Completed}$ \Comment{output-buffer ownership}
    \State $n_1\gets\min(B_r,N)$.
    \State $q[0]\gets\Call{DmaRead}{Q_{h,1,\rm DRAM},Q_{16}[0],2n_1d}$.
    \For{$i=1,\ldots,T_r$}
        \State $b\gets(i-1)\bmod2$; $r_0\gets(i-1)B_r$; $n_i\gets\min(B_r,N-r_0)$.
        \State $\Call{WaitAll}{q[b],w[b]}$ \Comment{input ready; output slot reusable}
        \If{$i<T_r$}
            \State $n_{i+1}\gets\min(B_r,N-iB_r)$.
            \State $q[1-b]\gets\Call{DmaRead}{Q_{h,i+1,\rm DRAM},Q_{16}[1-b],2n_{i+1}d}$.
            \Comment{prefetch while current tile computes}
        \EndIf
        \State $(Q_8,s_Q,U,m,\ell)\gets\VSync(\textsc{InitQuery},Q_{16}[b],n_i)$.
        \Statex \hspace{\algorithmicindent}\hspace{\algorithmicindent}\textit{InitQuery decodes and quantizes valid rows, zero-pads Q, then sets $U=0$, $m=-\infty$, $\ell=0$.}
        \State $J_i\gets\left\lceil\min(iB_r,N)/B_c\right\rceil$.
        \For{$j=1,\ldots,J_i$}
            \State $k_0\gets(j-1)B_c$.
            \State $S_{32}\gets\SafeDot(Q_8,K^{\rm dig}_j,B_c,d)$.
            \Comment{integer $Q_8K_{8,j}^{\top}$}
            \State $A\gets(S_{32},s_Q,s_K,m,\ell,r_0,k_0,n_i,N)$.
            \State $(P_8,\alpha,m,\ell)\gets\VSync(\textsc{SoftmaxPack},A)$.
            \State $C_{32}\gets\SafeDot(P_8,V^{\rm dig}_j,d,B_c)$.
            \Comment{integer $P_8V_{8,j}$; V is physically transposed}
            \State $U\gets\VSync(\textsc{Merge},U,\alpha,C_{32},s_V)$.
        \EndFor
        \State $O_{16}[b]\gets\VSync(\textsc{Finish},U,\ell,n_i,s_{O,h})$.
        \State $w[b]\gets\Call{DmaWrite}{O_{16}[b],O_{h,i,\rm DRAM},2n_id}$.
        \Comment{retain buffer until $w[b]$ completes}
    \EndFor
    \State $\Call{WaitAll}{w[0],w[1]}$.
\EndFor
\State \Return completed outputs for these four query heads.
\end{algorithmic}
\end{algorithm}

\textbf{What overlaps.} The next query tile can arrive by DMA while the current
tile computes. An output DMA can overlap later computation when dependencies
permit. Each DMA submission observes the single engine's acceptance rules;
simultaneous read/write service is not assumed. The two arrays run independent
microtiles in Algorithm 3. The inner softmax and output updates are explicitly
ordered, so neither array reads probabilities before the vector CPU finishes.

\textbf{Causal indexing.} All positions passed to vector routines are zero-based.
For row $r$ and column $c$, the query and key positions are $r_0+r$ and $k_0+c$.
Blocks with no causally visible key are skipped by $J_i$. Partially visible
blocks still require elementwise masking. The final normalization uses the
current $U,\ell$, which are precisely the state after block $J_i$.

\textbf{No direct compute access to DRAM.} Only DMA touches accelerator DRAM.
Array operand reads and result writes occur through SRAM. SRAM buffers for
array partials are retained until vector reconstruction has consumed them.

\clearpage
\begin{algorithm}[H]
\caption{Vector-CPU routines: online softmax, accumulation, and final quantization}
\label{alg:vector}
\small
\begin{algorithmic}[1]
\Function{SoftmaxPack}{$S_{32},s_Q,s_K,m,\ell,r_0,k_0,n_i,N$}
    \State $S\gets\F(S_{32})\,(s_Qs_K/\sqrt{d})$.
    \For{$r=0,\ldots,B_r-1$}
        \For{$c=0,\ldots,B_c-1$}
            \If{$r\ge n_i$ \textbf{or} $k_0+c\ge N$ \textbf{or} $k_0+c>r_0+r$}
                \State $S[r,c]\gets-\infty$.
            \EndIf
        \EndFor
        \State $m_{\rm new}\gets\max(m[r],\max_c S[r,c])$.
        \If{$m_{\rm new}=-\infty$}
            \State $\alpha[r]\gets0$; $P[r,:]\gets0$.
            \Comment{avoid $-\infty-(-\infty)$}
        \Else
            \If{$m[r]=-\infty$}
                \State $\alpha[r]\gets0$.
            \Else
                \State $\alpha[r]\gets\exp(m[r]-m_{\rm new})$.
            \EndIf
            \State $P[r,:]\gets\exp(S[r,:]-m_{\rm new})$.
        \EndIf
        \State $\ell[r]\gets\alpha[r]\ell[r]+\sum_cP[r,c]$; $m[r]\gets m_{\rm new}$.
        \State $P_8[r,:]\gets\I{8}\!\left(\clip(\round(127P[r,:]),0,127)\right)$.
    \EndFor
    \State \Return $P_8,\alpha,m,\ell$.
\EndFunction
\Statex
\Function{Merge}{$U,\alpha,C_{32},s_V$}
    \For{$r=0,\ldots,B_r-1$}
        \State $U[r,:]\gets\alpha[r]U[r,:]+(s_V/127)\F(C_{32}[r,:])$.
    \EndFor
    \State \Return $U$.
\EndFunction
\Statex
\Function{Finish}{$U,\ell,n_i,s_O$}
    \For{$r=0,\ldots,n_i-1$}
        \If{$\ell[r]>0$}
            \State $Y[r,:]\gets U[r,:]/\ell[r]$.
        \Else
            \State $Y[r,:]\gets0$.
        \EndIf
        \State $O_{16}[r,:]\gets\I{16}\!\left(\clip(\round(Y[r,:]/s_O),-32768,32767)\right)$.
    \EndFor
    \State \Return $O_{16}$.
\EndFunction
\end{algorithmic}
\end{algorithm}

\textbf{The essential correction.} Let $\alpha_r=\exp(m_r^{\rm old}-m_r^{\rm new})$.
The old output is \emph{multiplied} by $\alpha_r$, and the new integer product is
dequantized using both probability and value scales:
\[
 U^{\rm new}=\operatorname{diag}(\alpha)U^{\rm old}
       +\frac{s_V}{127}C_{32},\qquad
 \ell^{\rm new}=\alpha\odot\ell^{\rm old}+\rsum(P).
\]
Only the final output is divided by $\ell$. FP32 reads/writes do not themselves
introduce rounding; FP32 arithmetic and the explicit quantization steps do.

\textbf{Accuracy contract.} The denominator uses the FP32 exponential weights;
the numerator uses their INT8 approximation. This is a mixed-precision variant,
so it is not bitwise equivalent to full-precision attention. Fixed probability
scale $1/127$ can discard small weights. Quantization accuracy should be checked
on representative inputs before replacing a 16-bit activation kernel.

\clearpage
\textbf{An overflow-safe array primitive.}
For an INT8 value $x$, form signed digits on the vector CPU, with absolute values
evaluated after widening:
\[
 x_0=\operatorname{sign}(x)(|x|\bmod16),\qquad
 x_1=\operatorname{sign}(x)\lfloor|x|/16\rfloor,
 \qquad x=x_0+16x_1.
\]
Digits are stored in INT8 slots. Every digit magnitude is at most 15, so a
128-term partial dot product satisfies $|C_{16}|\le128\cdot15^2=28800<32768$.
With at most 256 original terms, reconstruction fits INT32.

\begin{algorithm}[H]
\caption{Control-CPU \SafeDot: $C=LW^{\top}$ using both systolic arrays}
\label{alg:dot}
\footnotesize
\begin{algorithmic}[1]
\Require $L\in\I{8}^{128\times R}$; cached right digits $W^{\rm dig}$ for
$W\in\I{8}^{n\times R}$; $n\in\{128,256\}$, $R\in\{128,256\}$.
\State $L^{\rm dig}\gets\VSync(\textsc{SplitAndPackLeft},L)$.
\Statex \hspace{\algorithmicindent}\textit{Pack contiguous microtiles, with each reduction segment at most 128 elements.}
\State $E\gets\varnothing$; assign a distinct SRAM destination to every raw partial.
\For{$t=0,128,\ldots,R-1$}
    \State $R_t\gets\min(128,R-t)$.
    \For{$a=0,1$}
        \For{$b=0,1$}
            \For{$c=0,16,\ldots,n-16$}
                \For{$u=0,1,2$}
                    \State $r\gets32u$; $\Call{WaitAccept}{\mathrm{SA}_{32}}$.
                    \State $(p_L,p_W,p_D)\gets\text{packed pointers for }(a,b,r,c,t,32)$.
                    \State $E\gets E\cup\{\Call{SaDot}{\mathrm{SA}_{32},p_L,p_W,p_D,R_t}\}$.
                    \If{$u<2$}
                        \State $r\gets96+16u$; $\Call{WaitAccept}{\mathrm{SA}_{16}}$.
                        \State $(p_L,p_W,p_D)\gets\text{packed pointers for }(a,b,r,c,t,16)$.
                        \State $E\gets E\cup\{\Call{SaDot}{\mathrm{SA}_{16},p_L,p_W,p_D,R_t}\}$.
                    \EndIf
                \EndFor
            \EndFor
        \EndFor
    \EndFor
\EndFor
\State $\Call{WaitAll}{E}$.
\State $C_{32}\gets\VSync(\textsc{Reconstruct},D)$.
\Statex \hspace{\algorithmicindent}\textit{Reconstruct initializes $C_{32}=0$ and adds $16^{a+b}\I{32}(D_{t,a,b,r,c})$ to each destination tile.}
\Statex \hspace{\algorithmicindent}\textit{Widen before multiplying; sum all digit pairs and reduction segments. No INT16 accumulation across commands.}
\State \Return $C_{32}$; release raw-partial storage only after reconstruction completes.
\end{algorithmic}
\end{algorithm}

\textbf{Layout and orientation.} A physical array operation computes
$D[r,c]=\sum_{k=0}^{R_t-1}L[r,k]W[c,k]$. The bracketed operands in Algorithm 3
denote pointers to already packed, contiguous blocks, not unsupported array
gathers or strided reads. \textsc{PrepareKV} caches K digits in 16-key by
128-feature blocks, and V digits in 16-feature by 128-key blocks, grouped by the
256-key tile. Thus the score call uses $W=K_{8,j}$, while the value call uses
$W=V_{8,j}^{\top}$ and computes $P_8(V_{8,j}^{\top})^{\top}=P_8V_{8,j}$.

\textbf{Scheduling and cost.} The 32-row array owns rows 0--95; the 16-row array
owns rows 96--127. This is a valid work split, not a measured optimum. Each
\textsc{SaDot} call still incurs 100 communication cycles and writes its INT16
tile through an 8 B/cycle output port. Distinct destinations allow the stated
array input/previous-output overlap when the array accepts the next operation.
Digit decomposition adds four physical products per reduction segment; these
commands, packing, and reconstruction must be included in any timing estimate.
Given the assignment's unbounded vector-compute assumption, a vector-only
product should also be compared using its SRAM traffic before claiming optimality.

\textbf{Public algorithm reference.} Tri Dao, \emph{FlashAttention-2}, Algorithm 1:
\url{https://tridao.me/publications/flash2/flash2.pdf}.
Hardware-specific scheduling and arithmetic handling above are adaptations to
the supplied specification.

\end{document}
